\section{Background and Related Work}\label{sec:background}

\subsection{Inverse problems on Markov chains}
\label{sec:inverse-mc}

A Markov chain is a stochastic process in which the next state depends
only on the current state. The \emph{forward} problem asks, given a
transition kernel $\Pchain$, what statistics (stationary distribution,
hitting times, expected occupation, etc.) does the chain produce? The
\emph{inverse} problem reverses this: given partial statistics observed
on the chain, recover the transition kernel itself. Inverse problems on
Markov chains arise in transport modelling \cite{ben-akiva1985},
biological systems \cite{noe2008}, and recommender systems
\cite{huang2018}, and are typically ill-posed under sparse observation
because many distinct kernels produce identical low-order statistics.
The usual remedy is a parametric restriction: rather than fit
$\Pchain \in \mathbb{R}^{n \times n}$ directly, one fits a parameter
vector $\theta \in \mathbb{R}^d$ with $d \ll n^2$ such that the
parametric form $\Pchain(\theta)$ is constrained to lie in a manifold of
plausible chains. The choice of parametric family is then load-bearing:
the wrong family caps recovery quality regardless of optimisation
effort.

\subsection{Route choice as a softmax-parametric Markov chain}
\label{sec:route-softmax}

In the route-choice setting the state space $\mathcal{X}$ is the set of
directed road edges and the kernel
$\pT(x'\mid x)$ is the probability that a vehicle currently on edge $x$
moves next to edge $x'$. Multinomial-logit route-choice models
\cite{ben-akiva1985, frejinger2007} parametrise the transition
probability as a softmax over the legal successors of $x$:
\begin{equation}
\pT(x'\mid x) \;=\;
\frac{\exp\bigl(s(x, x')\bigr)}{\sum_{x'' \in \mathrm{adj}(x)} \exp\bigl(s(x, x'')\bigr)},
\label{eq:softmax}
\end{equation}
where $s(x, x')$ is a real-valued score, typically a linear combination
of features of the edge $(x, x')$ and of the destination state $x'$.
This formulation has two important properties for inverse problems.
First, it is parameter-efficient: one global feature shared across
many states adds one parameter, regardless of network size. Second, its
gradient is analytic in $s$, which combined with the chain-rule
propagation through stationary computations and hitting-rate solves makes
$\nabla_\theta \log \pi$ and $\nabla_\theta \log \gobs$ tractable to
derive in closed form.

\subsection{The Morimura inverse-MC framework}
\label{sec:morimura}

\citet{morimura2013} apply the softmax form of \cref{eq:softmax} to the
partial-observation regime in which the modeller observes only a
small subset $\Xo \subset \mathcal{X}$. They consider two empirical
observation types:
\begin{itemize}
\item \textbf{Empirical stationary distribution $\fobs$ at $\Xo$.}
A count of vehicle passages at each station; effectively a sample of
$\pi(x)$ for $x \in \Xo$.
\item \textbf{Empirical hitting rates $\gobs(x, x')$ on $\Xo \times \Xo$.}
The fraction of vehicles observed at $x$ that subsequently visit $x'$
discounted by their travel time; effectively a sample of the
$\beta$-discounted first-passage probability.
\end{itemize}
They define a regularised loss that mixes a log-ratio match on $\fobs$
(via the stationary distribution of the parametric chain) and a log-error
match on $\gobs$ (via the hitting-rate matrix of the parametric chain),
weighted by a mixing parameter $\gamma$. Their key contribution is the
analytic derivation of the gradient of this loss with respect to
$\theta$, enabling optimisation via L-BFGS-B or, alternatively, the
natural-gradient descent the paper introduces in $\S 4.1$ as the
optimiser of choice when the loss landscape is ill-conditioned.

The paper validates the framework synthetically on a $n=100$ network
(their $\S 6.1$) and applies it to a $1497$-link Nairobi minibus
network as a real-world demonstration ($\S 6.2$). Both validations report
favourable RMAE on the recovered stationary distribution. Two limitations
of this validation are directly relevant to our work. First, the
synthetic experiment uses analytical (i.e.\ noise-free) hitting rates
$\gobs$ computed from the true chain; on noisy real-world data the
empirical $\gobs$ is overwhelmingly sparse, an effect documented in
\cref{sec:phase4-fg-infeasible}. Second, the real-world Nairobi case is
not framed as a regime-detection problem: the authors do not split the
data into two coexisting routing populations and ask whether the chain
recovery distinguishes them.

\subsection{Re-routing detection and external traffic influence}
\label{sec:rerouting-prior}

Detecting externally-influenced rerouting is a relatively recent
research direction. \citet{gu2018} proposed a junction-level
anomaly-detection approach in which deviations from a learned
turning-distribution baseline serve as evidence of unusual routing
behaviour; their setup is restricted to snapshot-level (population-level)
detection rather than per-trajectory. Active-routing-game analyses
\cite{thai2016, krichene2018} model the equilibrium effects of
sat-navigation adoption on aggregate network performance, but assume a
known navigation rule and focus on welfare bounds rather than detection.
Closest in spirit to this work is the line of route-choice research
that combines partially observed counts with route-fragment
likelihoods \cite{flotterod2017}: this addresses estimation under
partial observation but does not separate populations.

To our knowledge no prior work combines the Morimura inverse-MC
framework with a regime-detection objective, evaluated across both
controlled simulation and real-world data with an OD-mix-controlled
contrast. That is the gap this thesis addresses.

\subsection{Microscopic traffic simulation: SUMO}
\label{sec:sumo}

Simulation of Urban MObility (SUMO) \cite{lopez2018} is the open-source
microscopic traffic simulator used in this thesis to construct the
controlled validation tier (\cref{sec:phase4-sumo}). Each vehicle is
simulated with realistic car-following and lane-change dynamics on a
network derived from OpenStreetMap. Crucially, SUMO exposes a
\code{--device.rerouting.probability} flag that, when set to $1.0$,
gives each vehicle access to an internal congestion-aware
sat-navigation device; when set to $0.0$, vehicles route via static
shortest-path Dijkstra. By running an identical vehicle population
twice with that flag toggled, we obtain a controlled labelled
comparison --- intrinsic-routed vs.\ sat-navigation-routed --- against
which the Morimura-based detector can be honestly benchmarked.

\subsection{Real-world traffic from Automatic Vehicle Identification}
\label{sec:avi-prior}

AVI systems (Automatic License Plate Recognition cameras at selected
junctions, optionally extended with Bluetooth or floating-car-data
fusion) have produced a small but growing collection of openly
released city-scale traffic datasets. Among recent releases:
the Padua AVI deposit \cite{cappi2026} provides $10$-minute aggregated
flow statistics and pairwise travel times over $40$ junctions but
explicitly strips per-vehicle trajectories for privacy; the
Wang \emph{et al.}\ Xuancheng dataset \cite{wang2023xuancheng}
similarly restricts raw trajectory access; the
\citet{ma2026xuancheng} Xuancheng dataset --- used in this thesis ---
releases joinable entry/exit tables with anonymised vehicle IDs from
which complete trajectories can be reconstructed, covering
$\sim 350{,}000$ trips/day across $\sim 1{,}760$ road segments for the
month of April~2023. The Ma~\emph{et al.}\ release is, to our knowledge,
the largest openly available city-scale AVI dataset with
per-vehicle granularity at the time of writing.
